\section{Find Peak Element}
\label{sec:bs-peak-1d}

\problemstatement{We are given an array $\textit{arr}$ of $n$ distinct
integers (no two adjacent elements are equal), where the array is
conceptually bordered by $-\infty$ on both ends (i.e.\ $\textit{arr}[-1] =
\textit{arr}[n] = -\infty$). An index $i$ is a \textbf{peak} if
$\textit{arr}[i]$ is strictly greater than both of its neighbors. We must
find the index of \textbf{any one} peak element, in $O(\log n)$ time.}

\begin{patternbox}
This problem has \textbf{no sortedness at all} -- the array can be
completely jumbled. This is the most important lesson in this section:
the keyword to look for is not ``sorted,'' but the guarantee of a
\textbf{boundary condition combined with a local comparison rule} that
forces a peak to exist and forces a specific monotonic structure on the
direction of ``ascent.'' Whenever you see ``array bordered by $-\infty$''
or an equivalent guarantee (e.g.\ ``the array is strictly increasing then
strictly decreasing''), suspect a peak-finding binary search even though no
sortedness is mentioned.
\end{patternbox}

\subsection*{Understanding the problem carefully}

Consider the midpoint $\textit{mid}$ and compare $\textit{arr}[\textit{mid}]$
to its right neighbor $\textit{arr}[\textit{mid}+1]$ (treating an
out-of-range neighbor as $-\infty$ per the problem's border convention).

\begin{itemize}
  \item If $\textit{arr}[\textit{mid}] < \textit{arr}[\textit{mid}+1]$:
        the sequence is ``climbing'' at \textit{mid}. Because the array is
        bordered by $-\infty$ on the right end too, \emph{somewhere} to the
        right of \textit{mid} the climbing must eventually stop and turn
        into a peak (it cannot climb forever, since it must come back down
        to $-\infty$ at the far border) -- so a peak is guaranteed to exist
        in $(\textit{mid}, \textit{hi}]$.
  \item If $\textit{arr}[\textit{mid}] > \textit{arr}[\textit{mid}+1]$:
        the sequence is ``descending'' at \textit{mid}, meaning
        \textit{mid} is locally higher than its right neighbor. By a
        symmetric argument applied leftward (the array is bordered by
        $-\infty$ on the left too), a peak is guaranteed to exist somewhere
        in $[\textit{lo}, \textit{mid}]$, \emph{including possibly
        \textit{mid} itself}.
\end{itemize}

\begin{ideabox}[Greedy Choice]
If $\textit{arr}[\textit{mid}] < \textit{arr}[\textit{mid}+1]$, discard the
left half including \textit{mid}, since a peak is guaranteed further
right: $\textit{lo} \gets \textit{mid}+1$. Otherwise, discard the right
half excluding \textit{mid}, since a peak is guaranteed at \textit{mid} or
further left: $\textit{hi} \gets \textit{mid}$.
\end{ideabox}

\subsection*{Why a peak is always guaranteed to exist in the kept half}

This is an existence argument, not a uniqueness or value-comparison
argument, which makes it a structurally different kind of binary search
correctness proof from the rest of this chapter. The guarantee rests
entirely on the $-\infty$ borders: if we always move \emph{toward}
increasing values, we are guaranteed to eventually hit a local maximum
before falling off either border, because the function value cannot
increase forever in a finite array, and any point where increase changes
to decrease (or where increase meets the $-\infty$ border with the array
ascending all the way to its last real element) is, by definition, a peak.
This is the same idea used to prove that a continuous function climbing
from $-\infty$ must have a maximum -- a discrete, array-based analogue.

\subsection*{Algorithm}

\begin{enumerate}
  \item Initialize $\textit{lo} \gets 0$, $\textit{hi} \gets n-1$.
  \item While $\textit{lo} < \textit{hi}$:
  \begin{enumerate}
    \item $\textit{mid} \gets \textit{lo}+(\textit{hi}-\textit{lo})/2$.
    \item If $\textit{arr}[\textit{mid}] < \textit{arr}[\textit{mid}+1]$:
          $\textit{lo} \gets \textit{mid}+1$.
    \item Else: $\textit{hi} \gets \textit{mid}$.
  \end{enumerate}
  \item Return \textit{lo} (equivalently \textit{hi}; they have converged).
\end{enumerate}

\begin{algorithm}[H]
\caption{Find Peak Element}
\begin{algorithmic}[1]
\Procedure{FindPeak}{\textit{arr}}
  \State $\textit{lo} \gets 0$, $\textit{hi} \gets n-1$
  \While{$\textit{lo} < \textit{hi}$}
    \State $\textit{mid} \gets \textit{lo}+(\textit{hi}-\textit{lo})/2$
    \If{$\textit{arr}[\textit{mid}] < \textit{arr}[\textit{mid}+1]$}
      \State $\textit{lo} \gets \textit{mid} + 1$
    \Else
      \State $\textit{hi} \gets \textit{mid}$
    \EndIf
  \EndWhile
  \State \Return \textit{lo}
\EndProcedure
\end{algorithmic}
\end{algorithm}

\subsection*{Dry run}

\dryrun{Take $\textit{arr} = [1, 2, 1, 3, 5, 6, 4]$.}

\begin{itemize}
  \item $\textit{lo}=0,\textit{hi}=6$: $\textit{mid}=3$,
        $\textit{arr}[3]=3 < \textit{arr}[4]=5$. $\textit{lo}=4$.
  \item $\textit{lo}=4,\textit{hi}=6$: $\textit{mid}=5$,
        $\textit{arr}[5]=6 < \textit{arr}[6]=4$? No. $\textit{hi}=5$.
  \item $\textit{lo}=4,\textit{hi}=5$: $\textit{mid}=4$,
        $\textit{arr}[4]=5 < \textit{arr}[5]=6$. $\textit{lo}=5$.
  \item $\textit{lo}=5=\textit{hi}$: loop ends.
\end{itemize}

The peak index is $5$, where $\textit{arr}[5]=6$, which is indeed greater
than both neighbors $\textit{arr}[4]=5$ and $\textit{arr}[6]=4$. (Index $1$,
where $\textit{arr}[1]=2$, is also a valid peak; the problem only requires
\emph{any} one peak, and this algorithm happens to converge to the one on
the side it explored.)

\subsection*{C++ implementation}

\begin{lstlisting}
#include <bits/stdc++.h>
using namespace std;

int findPeak(vector<int>& arr) {
    int n = arr.size();
    int lo = 0, hi = n - 1;

    while (lo < hi) {
        int mid = lo + (hi - lo) / 2;
        if (arr[mid] < arr[mid + 1]) {
            lo = mid + 1;
        } else {
            hi = mid;
        }
    }
    return lo;
}

int main() {
    vector<int> arr = {1, 2, 1, 3, 5, 6, 4};
    cout << "Peak index: " << findPeak(arr) << endl;
    return 0;
}
\end{lstlisting}

\begin{complexitybox}
\textbf{Time Complexity.} The search window halves each iteration, giving
\[
  O(\log n).
\]
\textbf{Space Complexity.} $O(1)$ auxiliary space.
\end{complexitybox}

\begin{takeawaybox}
Binary search does not require sortedness -- it requires a monotonic,
provable rule for discarding half the search space. Here, the rule comes
from a local ascent/descent comparison plus a boundary guarantee, not from
any global ordering of the array's values. Whenever you see ``find a local
maximum/minimum,'' check whether boundary conditions and a single local
comparison can replace the sortedness argument used elsewhere in this
chapter.
\end{takeawaybox}
